\documentclass[a4paper,11pt]{article}
\usepackage[T2A]{fontenc}
\usepackage[utf8]{inputenc}
\usepackage[russian]{babel}
\usepackage{amsmath,amssymb,mathtools}
\usepackage{geometry}
\usepackage{microtype}
\usepackage{booktabs,tabularx,array}
\usepackage{enumitem}
\usepackage{xcolor}
\usepackage{tikz}
\usepackage{hyperref}
\usepackage{fancyhdr}
\usepackage{parskip}
\geometry{left=22mm,right=22mm,top=23mm,bottom=26mm}
\definecolor{accent}{HTML}{007D80}
\definecolor{accentdark}{HTML}{00575C}
\definecolor{ink}{HTML}{242D34}
\definecolor{softgray}{HTML}{68737C}
\definecolor{warn}{HTML}{AE3540}
\hypersetup{colorlinks=true,linkcolor=accentdark,urlcolor=accentdark,pdftitle={Ревью домашней работы — 15 ноября 2026 года},pdfauthor={Лёвин Артём Александрович}}
\setlength{\parindent}{0pt}
\setlength{\parskip}{4pt}
\setlist[itemize]{topsep=1pt,itemsep=1pt,leftmargin=6mm}
\pagestyle{fancy}
\fancyhf{}
\fancyhead[L]{\footnotesize\color{softgray} Математика / комбинаторика}
\fancyhead[R]{\footnotesize\color{softgray} 15.11.26}
\fancyfoot[C]{\parbox{\textwidth}{\centering\footnotesize\color{softgray}Лёвин Артём Александрович эксклюзивно для Марка Гукина\quad\textbullet\quad\thepage}}
\renewcommand{\headrulewidth}{0.3pt}
\renewcommand{\footrulewidth}{0pt}
\newcommand{\smallheading}[1]{\par\vspace{3pt}{\color{accentdark}\bfseries #1}\par\vspace{-2pt}}
\newcommand{\markerror}{\textcolor{warn}{\bfseries Ошибка:}}
\newcommand{\markidea}{\textcolor{blue!60!black}{\bfseries Ключевая идея:}}
\newcommand{\markresult}{\textcolor{teal!60!black}{\bfseries Итог:}}
\begin{document}
\thispagestyle{empty}
\begin{center}
\vspace*{37mm}
{\large\color{accentdark} МАТЕМАТИКА}\par
\vspace{18mm}
{\LARGE\bfseries\color{ink} Ревью домашней работы}\par
\vspace{7mm}
{\large\color{accentdark} Комбинаторика: принцип умножения,\par размещения и сочетания}\par
\vspace{30mm}
{\large Марк Гукин}\par
\vspace{16mm}
{\large 15 ноября 2026 года}\par
\vspace{8mm}
{\large Лёвин Артём Александрович}\par
\end{center}
\vfill
\begin{center}
\begin{tikzpicture}[x=1cm,y=1cm]
\draw[accent!70,line width=0.85pt] (-6,0) .. controls (-3,0.65) and (-1,-0.45) .. (1,0.08)
  .. controls (3,0.6) and (5,-0.26) .. (6,0);
\end{tikzpicture}
\end{center}
\clearpage

\section*{Сводная таблица}
В таблицу включены задания, решение которых осталось незавершённым. Перейдите по номеру, чтобы прочитать подробный разбор.
\par\vspace{7pt}
\renewcommand{\arraystretch}{1.5}
\small
\noindent\begin{tabularx}{\linewidth}{@{}p{13mm}p{26mm}X p{24mm}p{25mm}@{}}
\toprule
\textbf{№} & \textbf{Тема} & \textbf{Что требуется завершить} & \textbf{Ваш ответ} & \textbf{Верный ответ}\\
\midrule
\hyperref[task:8b]{\textbf{8(б)}} & Признак делимости на $4$ & Учесть все допустимые двузначные окончания и перестановки остальных цифр. & Не указан & $192$\\
\hyperref[task:10]{\textbf{10(а, б)}} & Выбор точек & Посчитать пары точек для отрезков и тройки точек для треугольников. & Не указаны & а) $66$; б) $220$\\
\bottomrule
\end{tabularx}
\normalsize
\vfill
{\footnotesize\color{softgray} При проверке рукописной записи учитываются только различимые действия. Неразборчивым символам не приписывается предполагаемый смысл.}
\clearpage

\thispagestyle{empty}
\vspace*{0.28\textheight}
\begin{center}
{\LARGE\bfseries\color{accentdark} Разбор ошибок}\par
\vspace{8mm}
{\large\color{softgray} Как восстановить подсчёт вариантов шаг за шагом}
\end{center}
\clearpage

\phantomsection\label{task:8b}
\section*{Задание 8(б). Шестизначные числа, кратные четырём}
\smallheading{Текст задания}
Используя каждую цифру из набора $1,2,3,4,5,6$ ровно один раз, найдите количество шестизначных чисел, кратных $4$.

\smallheading{Ваше решение и ответ}
В записи различимы отдельные множители $4,3,2,1,5$, далее стоит знак вопроса. Смысл всей цепочки множителей по фотографии определить надёжно нельзя. Окончательный числовой ответ отсутствует.

\smallheading{Где именно остаётся пробел}
\markerror\ Подсчёт вариантов не завершён. Из записи нельзя установить, что перечислены все допустимые окончания числа и что для каждого окончания правильно подсчитаны перестановки оставшихся цифр.

Последовательность $4\cdot3\cdot2\cdot1$ могла бы использоваться при расстановке четырёх оставшихся цифр. Этот фрагмент полезен, если именно такая роль множителей предполагалась, однако сам по себе он не определяет число шестизначных чисел, кратных четырём.

\smallheading{Почему этого недостаточно}
Признак делимости на $4$ относится к \emph{двум последним цифрам}. Например, число с окончанием $12$ делится на $4$, а с окончанием $14$ не делится. Поэтому сначала нужно определить все допустимые \emph{двузначные окончания}, затем посчитать варианты для четырёх первых позиций.

\smallheading{\markidea\ От чего зависит количество вариантов}
Есть два независимых выбора: допустимое окончание и порядок оставшихся четырёх цифр. Общее количество получается по правилу умножения.

\smallheading{Полное правильное решение}
\textbf{Шаг 1.} Перечислим двузначные числа, составленные из разных цифр набора $1,2,3,4,5,6$ и кратные $4$:
\[
12,\quad 16,\quad 24,\quad 32,\quad 36,\quad 52,\quad 56,\quad 64.
\]
Получили $8$ допустимых окончаний.

\textbf{Шаг 2.} После выбора окончания остаются четыре разные цифры. Их можно поставить на первые четыре позиции $4\cdot3\cdot2\cdot1=4!=24$ способами. Нуля в наборе нет, поэтому каждое размещение даёт шестизначное число.

\textbf{Шаг 3.} Для каждого из восьми окончаний возможны $24$ расстановки первых четырёх цифр:
\[
8\cdot4!=8\cdot24=\boxed{192}.
\]

\smallheading{Проверка}
Можно сгруппировать окончания по последней цифре:
\begin{itemize}
  \item последняя цифра $2$: окончания $12,32,52$ — три;
  \item последняя цифра $4$: окончания $24,64$ — два;
  \item последняя цифра $6$: окончания $16,36,56$ — три.
\end{itemize}
Итого $3+2+3=8$, что подтверждает подсчёт.

\markresult\ Таких шестизначных чисел $\boxed{192}$.

\smallheading{Задача на закрепление}
Из цифр $1,2,3,4,5$ составляют пятизначные числа, используя каждую цифру ровно один раз. Сколько полученных чисел кратны $4$?
\clearpage

\phantomsection\label{task:10}
\section*{Задание 10. Отрезки и треугольники на плоскости}
\smallheading{Текст задания}
На плоскости отмечены $12$ точек; никакие три из них не лежат на одной прямой. Сколько можно построить: а) отрезков с концами в этих точках; б) треугольников с вершинами в этих точках?

\smallheading{Ваше решение и ответ}
В подпункте «а» обозначены места для вычисления и стоит знак вопроса; законченной арифметической записи нет. Для подпункта «б» решение и ответ на фотографии не представлены.

\smallheading{Где именно остаётся пробел}
\markerror\ Не выполнен подсчёт количества способов выбрать две и три точки из двенадцати. Последнего законченного математического шага, от которого можно было бы продолжить вычисление, в записи нет.

\smallheading{Почему важно учитывать порядок}
Отрезок с концами в точках $A$ и $B$ совпадает с отрезком с концами в точках $B$ и $A$. Точно так же треугольник с вершинами $A,B,C$ остаётся тем же треугольником при любой перестановке названий вершин. Значит, при последовательном выборе точек каждый геометрический объект подсчитывается несколько раз. Эти повторы нужно исключить.

\smallheading{\markidea\ Как устроен выбор}
Для отрезка требуется выбрать \textbf{две} различные точки без учёта порядка. Для треугольника требуется выбрать \textbf{три} различные точки без учёта порядка. Условие об отсутствии трёх точек на одной прямой гарантирует, что каждая тройка действительно задаёт треугольник.

\smallheading{Полное правильное решение: а) отрезки}
\textbf{Шаг 1.} Первую точку можно выбрать $12$ способами. Вторую — $11$ способами, поскольку точки должны различаться. Значит, получили $12\cdot11=132$ \emph{упорядоченные} пары.

\textbf{Шаг 2.} Каждому отрезку соответствуют ровно две такие пары: $(A,B)$ и $(B,A)$. Разделим результат на $2$:
\[
\frac{12\cdot11}{2}=\boxed{66}.
\]
\textbf{Проверка.} Если упорядочить точки, можно посчитать отрезки так: из первой точки к $11$ последующим, из второй — к $10$, и так далее. Получаем $11+10+\ldots+1=66$.

\smallheading{Полное правильное решение: б) треугольники}
\textbf{Шаг 1.} Три различные точки в определённом порядке можно выбрать $12\cdot11\cdot10$ способами.

\textbf{Шаг 2.} Одну и ту же тройку вершин можно записать в $3\cdot2\cdot1=6$ различных порядках. Каждая тройка даёт один треугольник. Поэтому:
\[
\frac{12\cdot11\cdot10}{3\cdot2\cdot1}=\frac{1320}{6}=\boxed{220}.
\]

\textbf{Проверка.} Каждый из $66$ отрезков можно дополнить третьей вершиной $10$ способами. При этом каждый треугольник будет получен трижды — по каждой из его сторон. Следовательно, $66\cdot10:3=220$.

\markresult\ Ответ: а) $\boxed{66}$ отрезков; б) $\boxed{220}$ треугольников.

\smallheading{Задача на закрепление}
На плоскости отмечены $9$ точек, никакие три из которых не лежат на одной прямой. Сколько можно построить: а) отрезков с концами в этих точках; б) треугольников с вершинами в этих точках?
\end{document}
